\documentclass[11pt]{article}

% ----------------------------------------------------------- essential packages
\usepackage[hmargin=.75in, vmargin=.75in]{geometry}
%\usepackage{amsmath}
%\usepackage{amsfonts}
%\usepackage{amsthm}
%\usepackage{amssymb}
%\usepackage{booktabs} % somewhat essential, I guess
%\usepackage{setspace} % for double spacing, etc.
%\usepackage{multicol} % multiple columns
%\usepackage{tabularx} % tabularx
\usepackage{enumitem} % for manually setting indent of list environments
\usepackage{textcomp} % for typing degree symbol \textdegree
\usepackage{xcolor}

% ------------------------------------------------------------------- font stuff
%\usepackage[T1]{fontenc}
%\usepackage[rm,sfdefault,condensed,scale=1]{roboto}

%\usepackage[sfdefault,scale=1.05]{AlegreyaSans}

\usepackage{sourcesanspro}
\renewcommand{\familydefault}{\sfdefault}

%\usepackage[sfdefault]{noto}
%\pdfmapfile{=noto.map}

% -------------------------------------------------------------- hyperlink stuff
\usepackage[hidelinks]{hyperref}
%\hypersetup{colorlinks=true,urlcolor=def}

% ------------------------------------------------------------------- tikz stuff
\usepackage{tikz}
\tikzset{%
  every picture/.style={yscale=-1}%
}
\usetikzlibrary{positioning}

% ----------------------------------------------------------------------- macros
\newcommand{\textsb}[1]{ % semi-bold
  {\fontseries{sb}\selectfont #1}
}

\newcommand{\sect}[2]{ % #1 = text, #2 = icon
  \vspace{.75em}
  {#2}
  {\sectsize \bfseries \scshape \MakeUppercase{#1}}
  \vspace{.75em}
%
  %{\color{hl}
    %\vspace{-1em}
    %\noindent\rule{\textwidth}{.5pt}
    %\vspace{-1em}
  %}
}

\newcommand{\thing}[5]{%
  %\vspace{.25em}
  \href{#3}{\color{\titlecolor} \titlesize \textbf{#1}}
  \hfill {#2}%\\
  %{\fontseries{sb}\selectfont
    %\href{#3}{\emph{#4}}
  %}
  {\fontseries{m}\selectfont \small
    \begin{itemize}
    {#5}
    \end{itemize}
  }
}

\newcommand{\thingnolink}[3]{%
  %\vspace{.25em}
  {\color{\titlecolor} \titlesize \textbf{#1}}
  \hfill {#2}%
  {\fontseries{m}\selectfont \small
    \begin{itemize}
    {#3}
    \end{itemize}
  }
}

% -------------------------------------------------------------- layout settings
\def\colheight{60.3em}

\def\leftcolwidth{.375\textwidth}
\def\rightcolwidth{.585\textwidth}

\newcommand{\debugbox}[1]{
  \setlength{\fboxsep}{0pt}
  %\fbox % uncomment this line to see bounding boxes of columns
  {#1}
}

\setlength{\parindent}{0pt}
%\setlength{\parskip}{1em}

% --------------------------------------------------------------- style settings
%\definecolor{hl}{rgb}{.1,.4,.95}
%\definecolor{hl}{rgb}{.95,.55,.15}
%\definecolor{hl}{rgb}{.95,.45,.05}

\def\shadeofgray{.25} % main text color
\definecolor{def}{rgb}{\shadeofgray,\shadeofgray,\shadeofgray}

\definecolor{hl}{rgb}{.22,.45,.70} % highlight color

\def\iconcolor{hl}
\def\iconscale{.024}%{.0288}
\def\iconbaseline{-14pt}%{-16pt}
\def\contacticonscale{.018}
\def\contacticonbaseline{-9pt}

\def\titlecolor{hl}

\newcommand{\sectsize}{\LARGE} % size of section names
\newcommand{\titlesize}{\large} % size of "thing" titles

% --------------------------------------------------------------- begin document
\begin{document}
%
\setlist[itemize]{
  nosep,       % remove margins above and below itemize environment
  leftmargin=* % set indentation of itemize environment
}
%
\thispagestyle{empty} % remove page number
%
\color{def} % use default color
%
% ----------------------------------------------------------------------- header
%
\textbf{\textsc{\Huge Chester Lian}}
\hfill
\hfill
\hfill
{\normalsize % this controls the size of all contact info
  \begin{tikzpicture}
    [fill=\iconcolor, scale=\contacticonscale, baseline=\contacticonbaseline]
    \path (0,0) -- (24,0) -- (24,24) -- (0,24) -- cycle;
    \fill (20,4) -- (4,4) .. controls (2.9,4) and (2.01,4.9) .. (2.01,6) --
    (2,18) .. controls (2,19.1) and (2.9,20) .. (4,20) -- (20,20) .. controls
    (21.1,20) and (22,19.1) .. (22,18) -- (22,6) .. controls (22,4.9) and
    (21.1,4) .. (20,4) -- cycle(20,8) -- (12,13) -- (4,8) -- (4,6) -- (12,11) --
    (20,6) -- (20,8) -- cycle;
  \end{tikzpicture}
  chester.e.lian@gmail.com
%
  \hfill
%
  \begin{tikzpicture}
    [fill=\iconcolor, scale=\contacticonscale, baseline=\contacticonbaseline]
    \path (0,0) -- (24,0) -- (24,24) -- (0,24) -- cycle;
    \fill (12,2) .. controls (6.4800,2) and (2,6.4800) .. (2,12) .. controls
    (2,17.5200) and (6.4800,22) .. (12,22) .. controls (17.5200,22) and
    (22,17.5200) .. (22,12) .. controls (22,6.4800) and (17.5200,2) .. (12,2) --
    cycle(11,19.9300) .. controls (7.0500,19.4400) and (4,16.0800) .. (4,12) ..
    controls (4,11.3800) and (4.0800,10.7900) .. (4.2100,10.2100) -- (9,15) --
    (9,16) .. controls (9,17.1000) and (9.9000,18) .. (11,18) -- (11,19.9300) --
    cycle(17.9000,17.3900) .. controls (17.6400,16.5800) and (16.9000,16) ..
    (16,16) -- (15,16) -- (15,13) .. controls (15,12.4500) and (14.5500,12) ..
    (14,12) -- (8,12) -- (8,10) -- (10,10) .. controls (10.5500,10) and
    (11,9.5500) .. (11,9) -- (11,7) -- (13,7) .. controls (14.1000,7) and
    (15,6.1000) .. (15,5) -- (15,4.5900) .. controls (17.9300,5.7800) and
    (20,8.6500) .. (20,12) .. controls (20,14.0800) and (19.2000,15.9700) ..
    (17.9000,17.3900) -- cycle;
    % below code is for a house icon
    %\path[fill] (10,20) -- (10,14) -- (14,14) -- (14,20) -- (19,20) -- (19,12)
    %-- (22,12) -- (12,3) -- (2,12) -- (5,12) -- (5,20) -- cycle;
    %\path (0,0) -- (24,0) -- (24,24) -- (0,24) -- cycle;
  \end{tikzpicture}
  need.to.fix.my.website
  %chesterlian.com
%
  \hfill
%
  \begin{tikzpicture}
    [fill=\iconcolor, scale=\contacticonscale, baseline=\contacticonbaseline]
    \path (0,0) -- (24,0) -- (24,24) -- (0,24) -- cycle;
    \fill (6.62,10.79) .. controls (8.06,13.62) and (10.38,15.93) ..
    (13.21,17.38) -- (15.41,15.18) .. controls (15.68,14.91) and (16.08,14.82)
    .. (16.43,14.94) .. controls (17.55,15.31) and (18.76,15.51) .. (20,15.51)
    .. controls (20.55,15.51) and (21,15.96) .. (21,16.51) -- (21,20) ..
    controls (21,20.55) and (20.55,21) .. (20,21) .. controls (10.61,21) and
    (3,13.39) .. (3,4) .. controls (3,3.45) and (3.45,3) .. (4,3) -- (7.5,3) ..
    controls (8.05,3) and (8.5,3.45) .. (8.5,4) .. controls (8.5,5.25) and
    (8.7,6.45) .. (9.07,7.57) .. controls (9.18,7.92) and (9.1,8.31) ..
    (8.82,8.59) -- (6.62,10.79) -- cycle;
  \end{tikzpicture}
  678-902-4787
}

{\color{hl}
  \vspace{-.5em}
  \rule{\textwidth}{1.25pt}
  \vspace{-2em}
}

% ------------------------------------------------------------------ left column
%
%\setlength{\fboxsep}{0pt}
\debugbox{
\begin{minipage}[t][\colheight]{\leftcolwidth}
\setlength{\parskip}{.25em}

  \sect{Education}{
    \begin{tikzpicture}
      [fill=\iconcolor, scale=\iconscale, baseline=\iconbaseline]
      \path (0,0) -- (24,0) -- (24,24) -- (0,24) -- cycle;
      \fill (5,13.18) -- (5,17.18) -- (12,21) -- (19,17.18) -- (19,13.18) --
      (12,17) -- (5,13.18) -- cycle(12,3) -- (1,9) -- (12,15) -- (21,10.09) --
      (21,17) -- (23,17) -- (23,9) -- (12,3) -- cycle;
    \end{tikzpicture}
  }

  \textbf{\titlesize Duke University}
  \hfill Durham, NC\\
  \textsb{Ph.D.}
  \hfill Expected May 2019

  {\small
    \begin{itemize}
      \item \textsb{Focus:} Mathematics (Algebraic Topology)
      \item \textsb{GPA:} 3.97
    \end{itemize}
  }

  \vfill

  \textbf{\titlesize University of Dayton}
  \hfill Dayton, OH\\
  \textsb{B.S.\ (Honors with Distinction)}
  \hfill 2013

  {\small
    \begin{itemize}
      \item \textsb{Majors:} Mathematics, Computer Science
      \item \textsb{Minor:} Philosophy
      \item \textsb{GPA:} 3.98 (summa cum laude)
      \item \textsb{Honors Thesis:} An Aronszajn Tree
      \item \textsb{Publication:} On the Decompositions of Complete Graphs into
      Cycles and Stars on the Same Number of Edges (with A. Abueida).
      \emph{Discussiones Mathematicae Graph Theory} 34.1 (2014), 113--125.
      \item \textsb{Putnam Competition:} 20 (2010), 21 (2011)
    \end{itemize}
  }

  \vfill\vfill
  \sect{Skills}{
    \begin{tikzpicture}
      [fill=\iconcolor, scale=\iconscale, baseline=\iconbaseline]
      \path (0,0) -- (24,0) -- (24,24) -- (0,24) -- cycle;
      \fill (22.7,19) -- (13.6,9.9) .. controls (14.5,7.6) and (14,4.9) ..
      (12.1,3) .. controls (10.1,1) and (7.1,0.6) .. (4.7,1.7) -- (9,6) -- (6,9)
      -- (1.6,4.7) .. controls (0.4,7.1) and (0.9,10.1) .. (2.9,12.1) ..
      controls (4.8,14) and (7.5,14.5) .. (9.8,13.6) -- (18.9,22.7) .. controls
      (19.3,23.1) and (19.9,23.1) .. (20.3,22.7) -- (22.6,20.4) .. controls
      (23.1,20) and (23.1,19.3) .. (22.7,19) -- cycle;
    \end{tikzpicture}
  }

  %\textsb{Programming:} JavaScript (CoffeeScript, Node.js), C++, Python, Perl

  %\vspace{.5em}

  %\textsb{Other:} SQL, HTML, CSS, \LaTeX, Mathematica, UNIX/Linux

  JavaScript (Node.js, CoffeeScript, jQuery), C++, Python, Perl, SQL, \LaTeX\
  (Ti\emph{k}Z), HTML, CSS, Mathematica, UNIX/Linux

  \vfill\vfill
  \sect{Work Experience}{
    \begin{tikzpicture}
      [fill=\iconcolor, scale=\iconscale, baseline=\iconbaseline]
      \path (0,0) -- (24,0) -- (24,24) -- (0,24) -- cycle;
      \fill (20,6) -- (16,6) -- (16,4) .. controls (16,2.89) and (15.11,2) ..
      (14,2) -- (10,2) .. controls (8.89,2) and (8,2.89) .. (8,4) -- (8,6) --
      (4,6) .. controls (2.89,6) and (2.01,6.89) .. (2.01,8) -- (2,19) ..
      controls (2,20.11) and (2.89,21) .. (4,21) -- (20,21) .. controls
      (21.11,21) and (22,20.11) .. (22,19) -- (22,8) .. controls (22,6.89) and
      (21.11,6) .. (20,6) -- cycle(14,6) -- (10,6) -- (10,4) -- (14,4) -- (14,6)
      -- cycle;
    \end{tikzpicture}
  }

  \textbf{\titlesize Course Coordinator}\\
  \textsb{Lab Calculus \& Functions II} \hfill Spring 2018

  {\small
    \begin{itemize}

      \item Schedule and lead weekly meetings with all course instructors and
      lab TAs

      \item Guide less-experienced instructors and TAs in lesson planning and
      lab management

      \item Review exams written by other instructors and provide feedback

    \end{itemize}
  }

  \vfill

  \textbf{\titlesize Instructor of Record}\\
  \textsb{Teaching College Mathematics} \hfill Fall 2016

  {\small
    \begin{itemize}

      \item Co-instructed with a full-time professor

      \item Trained first-year graduate students in running and teaching lab
      calculus courses

    \end{itemize}
  }

  \vfill

  \textbf{\titlesize Instructor of Record}\\
  \textsb{Various calculus courses} \hfill 5 semesters

  {\small
    \begin{itemize}

      \item \textsb{Recognition:} Smith Award for Teaching Excellence (lump sum
      cash prize), academic years 2015--2016 and 2016--2017

    \end{itemize}
  }

\end{minipage}
} % end of debugbox
%
\hfill
%
% ----------------------------------------------------------------- right column
%
%\setlength{\fboxsep}{0pt}
\debugbox{
\begin{minipage}[t][\colheight]{\rightcolwidth}
\setlength{\parskip}{.25em}

  \sect{Independent Projects}{
    \begin{tikzpicture}
      [fill=\iconcolor, scale=\iconscale, baseline=\iconbaseline]

      \path (0,0) -- (24,0) -- (24,24) -- (0,24) -- cycle;

      \fill (20.5,11) -- (19,11) -- (19,7) .. controls (19,5.9) and (18.1,5) ..
      (17,5) -- (13,5) -- (13,3.5) .. controls (13,2.12) and (11.88,1) ..
      (10.5,1) .. controls (9.12,1) and (8,2.12) .. (8,3.5) -- (8,5) -- (4,5) ..
      controls (2.9,5) and (2.01,5.9) .. (2.01,7) -- (2.01,10.8) -- (3.5,10.8)
      .. controls (4.99,10.8) and (6.2,12.01) .. (6.2,13.5) .. controls
      (6.2,14.99) and (4.99,16.2) .. (3.5,16.2) -- (2,16.2) -- (2,20) ..
      controls (2,21.1) and (2.9,22) .. (4,22) -- (7.8,22) -- (7.8,20.5) ..
      controls (7.8,19.01) and (9.01,17.8) .. (10.5,17.8) .. controls
      (11.99,17.8) and (13.2,19.01) .. (13.2,20.5) -- (13.2,22) -- (17,22) ..
      controls (18.1,22) and (19,21.1) .. (19,20) -- (19,16) -- (20.5,16) ..
      controls (21.88,16) and (23,14.88) .. (23,13.5) .. controls (23,12.12) and
      (21.88,11) .. (20.5,11) -- cycle;

    \end{tikzpicture}
  }

  \thing
    {Calculus Course Website}
    {Spring 2017}
    {http://math.duke.edu/~chester/teaching}
    {math.duke.edu/\textasciitilde chester/teaching}
    {
      \item Built course website for a calculus class

      \item Developed familiarity with AngularJS components and AngularJS
      Material (UI component framework)
    }

  \vfill

  \thing
    {360\textdegree{} 3D Rubik's Cube Visualization}
    {Summer 2015}
    {http://www.chesterlian.com/seenjs/demo.html}
    {chesterlian.com/seenjs/demo.html}
    {
      \item Implemented a 3D version of VisualCube
      %(\href{http://cube.crider.co.uk/visualcube.php}
      %{\emph{cube.crider.co.uk/visualcube.php}}),
      (PHP script for generating 2D PNGs of partial cube states) with basic
      features

      \item Rendered on HTML5 Canvas using seen.js (3D library)
    }

  \vfill

  \thing
    {Visualization of Calculus Concepts}
    {2014--2016}
    {http://math.duke.edu/~chester}
    {math.duke.edu/\textasciitilde chester}
    {
      \item Created interactive web apps to illustrate (pre-)calculus concepts
      to students, e.g., derivatives, phase planes, Euler's method
      \item Used jQuery and Flot (JavaScript plotting library)
    }

  \vfill

  \thing
    {PengyTimer}
    {Summer 2010}
    {http://www.chesterlian.com}
    {chesterlian.com}
    {
      \item Designed and developed fully keyboard-controlled Rubik's cube timer
      with common functionalities and minimal visual clutter (inspired by Vim)
    }

  \vfill\vfill
  \sect{School Projects}{
    \begin{tikzpicture}
      [fill=\iconcolor, scale=\iconscale, baseline=\iconbaseline]
      \path (0,0) -- (24,0) -- (24,24) -- (0,24) -- cycle; % bounding square
      \fill (19,3) -- (14.82,3) .. controls (14.4,1.84) and (13.3,1) .. (12,1)
      .. controls (10.7,1) and (9.6,1.84) .. (9.18,3) -- (5,3) .. controls
      (3.9,3) and (3,3.9) .. (3,5) -- (3,19) .. controls (3,20.1) and (3.9,21)
      .. (5,21) -- (19,21) .. controls (20.1,21) and (21,20.1) .. (21,19) --
      (21,5) .. controls (21,3.9) and (20.1,3) .. (19,3) -- cycle(12,3) ..
      controls (12.55,3) and (13,3.45) .. (13,4) .. controls (13,4.55) and
      (12.55,5) .. (12,5) .. controls (11.45,5) and (11,4.55) .. (11,4) ..
      controls (11,3.45) and (11.45,3) .. (12,3) -- cycle(14,17) -- (7,17) --
      (7,15) -- (14,15) -- (14,17) -- cycle(17,13) -- (7,13) -- (7,11) --
      (17,11) -- (17,13) -- cycle(17,9) -- (7,9) -- (7,7) -- (17,7) -- (17,9) --
      cycle;
    \end{tikzpicture}
  }

  \thingnolink
    {Ph.D.\ Dissertation}
    {Summer 2017}
    {
      \item Improved runtime of a search algorithm from several days to less
      than a minute

      \item more info
    }

  \vfill

  \thingnolink
    {``Facebook'' Project}
    {Spring 2012}
    {
      \item Implemented rudimentary version of Facebook using MySQL, PHP,
      JavaScript for Database Management Systems class with basic functionality
      (add/remove friends, wall posts, comments)

      %\item more info
    }

  \vfill

  \thing
    {University of Dayton Math Club Website}
    {2011--2013}
    {http://academic.udayton.edu/mathclub}
    {academic.udayton.edu/mathclub}
    {
      \item Overhauled content and design of defunct math club website
    }

  \vfill\vfill
  \sect{Activities}{
    \begin{tikzpicture}
      %[scale=.384, baseline=-6.2pt]
      [scale=.321, baseline=-5.8pt]
      \path (-.9,-.9) rectangle (.9,.9); % bounding square
      \definecolor{customred}{rgb}{.95,.2,.2} % some kind of soft red
      \begin{scope}[fill=customred]
        \fill (0.2405,-0.0231) -- (0.3737,-0.1593) --
          (0.3616,0.0738) -- (0.2334,0.2169) -- cycle;
        \fill (0.4152,-0.2034) -- (0.5347,-0.3256) --
          (0.5185,-0.0994) -- (0.4030,0.0296) -- cycle;
        \fill (0.5720,-0.3654) -- (0.6798,-0.4755) --
          (0.6604,-0.2560) -- (0.5558,-0.1392) -- cycle;
        \fill (0.2303,0.2927) -- (0.3585,0.1496) --
          (0.3473,0.3643) -- (0.2238,0.5130) -- cycle;
        \fill (0.3985,0.1033) -- (0.5140,-0.0257) --
          (0.4990,0.1836) -- (0.3873,0.3181) -- cycle;
        \fill (0.5501,-0.0675) -- (0.6547,-0.1843) --
          (0.6367,0.0196) -- (0.5352,0.1418) -- cycle;
        \fill (0.2210,0.5828) -- (0.3445,0.4341) --
          (0.3342,0.6328) -- (0.2150,0.7856) -- cycle;
        \fill (0.3831,0.3862) -- (0.4948,0.2518) --
          (0.4809,0.4461) -- (0.3728,0.5849) -- cycle;
        \fill (0.5299,0.2083) -- (0.6314,0.0862) --
          (0.6146,0.2761) -- (0.5160,0.4026) -- cycle;
      \end{scope}
      \definecolor{customyellow}{RGB}{255,215,0} % this is just "gold" in CSS
      \begin{scope}[fill=customyellow]
        \fill (-0.1592,-0.7281) -- (0.0476,-0.6754) --
          (-0.0760,-0.5770) -- (-0.2892,-0.6349) -- cycle;
        \fill (0.1196,-0.6564) -- (0.3378,-0.6007) --
          (0.2216,-0.4968) -- (-0.0039,-0.5580) -- cycle;
        \fill (0.4140,-0.5807) -- (0.6446,-0.5219) --
          (0.5368,-0.4118) -- (0.2978,-0.4767) -- cycle;
        \fill (-0.3336,-0.6014) -- (-0.1204,-0.5434) --
          (-0.2566,-0.4350) -- (-0.4765,-0.4989) -- cycle;
        \fill (-0.0460,-0.5225) -- (0.1796,-0.4612) --
          (0.0511,-0.3463) -- (-0.1822,-0.4141) -- cycle;
        \fill (0.2585,-0.4390) -- (0.4975,-0.3741) --
          (0.3780,-0.2520) -- (0.1301,-0.3240) -- cycle;
        \fill (-0.5255,-0.4618) -- (-0.3056,-0.3979) --
          (-0.4565,-0.2778) -- (-0.6834,-0.3487) -- cycle;
        \fill (-0.2287,-0.3747) -- (0.0046,-0.3069) --
          (-0.1382,-0.1791) -- (-0.3796,-0.2546) -- cycle;
        \fill (0.0864,-0.2822) -- (0.3344,-0.2101) --
          (0.2012,-0.0739) -- (-0.0564,-0.1544) -- cycle;
      \end{scope}
      \definecolor{customblue}{rgb}{.22,.45,.7} % same as hl color
      \begin{scope}[fill=customblue]
        \fill (-0.7074,-0.2912) -- (-0.4806,-0.2203) --
          (-0.4671,0.0117) -- (-0.6855,-0.0632) -- cycle;
        \fill (-0.4041,-0.1952) -- (-0.1627,-0.1198) --
          (-0.1588,0.1162) -- (-0.3907,0.0367) -- cycle;
        \fill (-0.0813,-0.0931) -- (0.1762,-0.0126) --
          (0.1691,0.2274) -- (-0.0775,0.1429) -- cycle;
        \fill (-0.6797,0.0105) -- (-0.4614,0.0853) --
          (-0.4489,0.2992) -- (-0.6594,0.2212) -- cycle;
        \fill (-0.3878,0.1116) -- (-0.1560,0.1910) --
          (-0.1525,0.4081) -- (-0.3754,0.3255) -- cycle;
        \fill (-0.0779,0.2189) -- (0.1686,0.3034) --
          (0.1621,0.5237) -- (-0.0744,0.4360) -- cycle;
        \fill (-0.6541,0.2894) -- (-0.4436,0.3674) --
          (-0.4321,0.5653) -- (-0.6353,0.4848) -- cycle;
        \fill (-0.3728,0.3946) -- (-0.1499,0.4772) --
          (-0.1466,0.6776) -- (-0.3613,0.5925) -- cycle;
        \fill (-0.0748,0.5060) -- (0.1617,0.5936) --
          (0.1557,0.7965) -- (-0.0716,0.7064) -- cycle;
      \end{scope}
    \end{tikzpicture}
  }

  \thingnolink
    {{\color{def} Rubik's Cube}}
    {}
    {
      \item \textsb{Official World Records:} Memorized and solved a 5x5x5
      Rubik's cube blindfolded in 15 minutes (July 2010), memorized and solved
      17 consecutive Rubik's cubes blindfolded in 53 minutes (June 2011)

      \item \textsb{TED Talk:} Invited to speak at TEDxBeaconStreet (November
      2014)

      \item \textsb{Media:} Guest on Anderson Cooper's \emph{Anderson LIVE}
      (September 2012)

      \item \textsb{Leadership:} (Co-)organized 7 World Cube Association
      competitions

      \item \textsb{Teamwork:} Staff member at 30+ competitions (World
      Championships 2013, US National Championships 2010, 2011, 2012, 2014,
      2015)
    }

\end{minipage}
} % end of debugbox

\end{document}
